% =====================================================================
\subsection{I moduli YANG}
\label{sec:yang018}
% =====================================================================

L'Annex~A (normativo) traduce il modello in tre moduli YANG~1.1, pubblicati
su ETSI Forge~\cite{qkd018yang}. Con i costrutti visti in
§\ref{sec:yang-anatomia} la loro struttura si legge rapidamente. I moduli
importano \texttt{etsi-qkd-sdn-node} ed \texttt{etsi-qkd-node-types} della
GS~QKD~015 per riusarne i tipi (identificativi, tipi di link, QoS delle
applicazioni): la 018 non reinventa il nodo, ne fornisce una vista
astratta. I tre moduli formano una catena: \emph{topology} (lo scheletro)
$\to$ \emph{inventory} (che cosa c'è e in che stato) $\to$
\emph{connectivity} (che cosa ci si fa).

\paragraph{\texttt{etsi-qkd-sdn-topology} (prefisso \texttt{qkdnw-tplg}).}
Descrive solo la struttura, cioè chi è collegato a chi:
\begin{lstlisting}
qkd_network
 +-- sdqkd_nodes
 |    +-- qkdn* [qkdn_id]                          nodi SD-QKD
 |         +-- qkd_interfaces/qkdi* [qkdi_id]      moduli QKD
 +-- qkd_links
      +-- qkd_phys_links/phys_link* [phys_link_id]  link fisici
      |    +-- link_type
      |    +-- phys_link_local_qkdn   (qkdn_id_ref, qkdi_id_ref)
      |    +-- phys_link_remote_qkdn  (qkdn_id_ref, qkdi_id_ref)
      +-- qkd_svc_links/svc_link* [svc_link_id]     service link
           +-- ... estremi locale/remoto come sopra
\end{lstlisting}
Rispetto alla 015 il punto di vista cambia: non c'è più un albero per
nodo, ma un unico albero per l'intera rete del dominio. In termini astratti
è un grafo $G=(V,E)$ con $V$ insieme dei nodi \texttt{qkdn} ed
$E = E_{\text{phys}}\cup E_{\text{svc}}$, dove i \texttt{phys\_link}
rappresentano adiacenze fisiche e i \texttt{svc\_link} relazioni logiche di
chiave, anche multi-hop. Gli estremi sono \texttt{leafref}, cioè
riferimenti validati a nodi e interfacce esistenti. Avendo il grafo $G_i$ di
ciascun dominio $i$, l'orchestrator compone una vista multi-dominio
\begin{equation}
  G = \Big(\bigcup_i V_i,\; \bigcup_i E_i \,\cup\, E_{\text{inter}}\Big),
\end{equation}
dove $E_{\text{inter}}$ sono le adiacenze tra domini, tipicamente
realizzate da nodi di bordo condivisi.

\paragraph{\texttt{etsi-qkd-sdn-inventory} (prefisso \texttt{qkdnw-invt}).}
Non crea un nuovo albero ma fa \texttt{augment} di quello di topologia,
aggiungendo foglie agli elementi già definiti (Tab.~\ref{tab:inventory}).
È un esempio da manuale dell'uso di \texttt{augment}: chi è interessato
solo alla struttura implementa la topologia, chi vuole anche lo stato
aggiunge l'inventario senza duplicare nulla.

\begin{table}[h]
\centering\small
\caption{Elementi aggiunti dal modulo inventory all'albero di topologia.}
\label{tab:inventory}
\begin{tabularx}{\textwidth}{@{}l L@{}}
\toprule
\textbf{Elemento} & \textbf{Foglie aggiunte} \\
\midrule
\texttt{qkdn} & \texttt{qkdn\_version}, \texttt{qkdn\_capabilities},
\texttt{qkdn\_location\_id}, \texttt{KME\_ip\_port} (IP e porta del KME, dove
i SAE prelevano le chiavi via QKD~014), \texttt{qkdn\_status} \\
\texttt{qkdi} & \texttt{role\_support}, \texttt{qkdi\_model},
\texttt{qkdi\_type}, \texttt{qkdi\_att\_point}, \texttt{qkdi\_status} \\
\texttt{phys\_link} & \texttt{enable} (configurabile), \texttt{wavelength},
\texttt{performance} (\texttt{skr}, \texttt{eskr},
\texttt{expected\_consumption} e la lista \texttt{phys\_link\_perf} con QBER
e SNR), \texttt{phys\_link\_status} \\
\texttt{svc\_link} & \texttt{bandwidth}, \texttt{performance},
\texttt{path\_restoration} (\texttt{type} \texttt{AUTOMATIC}/\texttt{MANUAL},
\texttt{hold\_off\_time} in ms), \texttt{svc\_link\_status} \\
\bottomrule
\end{tabularx}
\end{table}

Il modulo definisce anche le notifiche che realizzano la parte reattiva
della §5.5: \texttt{sdqkdn\_new}, \texttt{sdqkdn\_interface\_new},
\texttt{sdqkdn\_phys\_link\_new}, \texttt{sdqkdn\_svc\_link\_new}, le
varianti \texttt{\_perf\_update} e \texttt{\_overloaded}, distinte per link
fisici e service link (ad esempio \texttt{sdqkdn\_phys\_link\_overloaded}),
e le generiche \texttt{alarm} ed \texttt{event}, che riportano stato,
messaggio e severità di nodi, interfacce e link. Le notifiche sullo stato del
nodo elencate nella §5.5 (\texttt{sdqkdn\_operating}, \texttt{\_down},
\texttt{\_failure}, \texttt{\_out}) non hanno una \texttt{notification}
dedicata nel modulo: lo stato viaggia nel campo \texttt{qkdn\_status} di
\texttt{alarm} ed \texttt{event}. Il campo \texttt{eskr} dei link è il dato
da cui l'orchestrator capisce se un cammino regge un nuovo servizio,
secondo la~\eqref{eq:admissible}.

\paragraph{\texttt{etsi-qkd-sdn-connectivity} (prefisso \texttt{qkdnw-conn}).}
È il modulo operativo: gestisce le applicazioni e i cammini che le servono.
\begin{lstlisting}
connectivity_service/qkd_service/qkd_app* [app_id]
 +-- (dati dell'applicazione da QKD 015: stato, tipo, QoS, statistiche)
 +-- app_mapped_svc_link
      +-- svc_link_and_qkdn_ids   svc_link_id_ref + estremi (nodo, interfaccia)
      +-- path_constraints
      |    +-- include_nodes* [qkdn_id_ref qkdi_id_ref]
      |    +-- exclude_nodes* [qkdn_id_ref qkdi_id_ref]
      +-- candidate_paths
      |    +-- num_of_candidate_paths
      |    +-- candidate_path* [path_id]   phys_links*  (path_id = rango)
      +-- deploy_path             path_id + phys_links* del cammino scelto
\end{lstlisting}
Qui compaiono le \texttt{rpc}, che ricalcano passo per passo il workflow di
provisioning:
\begin{description}[font=\normalfont\ttfamily\footnotesize,style=nextline,leftmargin=1.2em]
  \item[app-registration] registra l'applicazione nel dominio;
  \item[create-service-link-id-for-app] crea il service link tra nodo
        locale e remoto;
  \item[create-candidate-paths-per-service-link-for-app] il controller
        calcola i cammini candidati rispettando QoS e vincoli
        \texttt{include\_nodes}/\texttt{exclude\_nodes}, espressi come
        coppie (nodo, interfaccia);
  \item[update-candidate-paths-per-service-link-for-app] ricalcolo con
        vincoli diversi;
  \item[deploy-a-candidate-path-per-service-link-for-app] l'orchestrator
        sceglie un \texttt{path\_id} e il controller lo attiva;
  \item[delete-deployed-path-per-service-link-for-app] rilascio del
        cammino.
\end{description}
Le notifiche del modulo sono \texttt{sdqkdn\_application\_new},
\texttt{sdqkdn\_application\_deployed},
\texttt{sdqkdn\_application\_qos\_update} e
\texttt{sdqkdn\_application\_disconnected}.

Il punto concettuale è la \textbf{divisione dei ruoli}: il controller
conosce il proprio dominio e \emph{propone} i cammini, l'orchestrator ha la
vista d'insieme e \emph{decide}. La Tab.~\ref{tab:moduli} riassume.

\begin{table}[h]
\centering\small
\caption{I tre moduli YANG della GS~QKD~018.}
\label{tab:moduli}
\begin{tabularx}{\textwidth}{@{}l L l L@{}}
\toprule
\textbf{Modulo} & \textbf{Domanda} & \textbf{Clausola} & \textbf{Dati} \\
\midrule
topology & Com'è fatta la rete? & §5.1 & struttura (grafo) \\
inventory & Quante risorse ci sono e in che stato? & §5.2, §5.5 &
stato, prestazioni, configurazione, notifiche \\
connectivity & Come servo un'applicazione end-to-end? & §5.3, §5.4, §6.2 &
applicazioni, RPC di calcolo e deploy dei cammini \\
\bottomrule
\end{tabularx}
\end{table}
